For a completed pair, the boost bivector $\sigma$ satisfies $\sigma^2 = 1$, generates the rotor flow with weights $\mp 1$ on the pair, has the occupation projectors as spectral projectors, acts in Fock representation through $\tfrac12\sigma = N - \tfrac12$, and is the unique one-parameter rotor flow of the atom \emph{up to reparametrisation of the flow parameter and orientation} (the bivector space is one-dimensional: the algebra fixes the orbits, not the speed, which is the normalisation the geometry supplies). The geometric normalisation converts rapidity to Schwarzschild time at the surface gravity $\kappa = 1/4M$. Continued to imaginary rapidity, the boost rotor becomes a unitary of the complexified atom whose conjugation action has period exactly $2\pi$; the continuation's generator $i\sigma$ lies outside the real algebra, so the cycle belongs to the clock's Euclidean continuation, not to any intrinsic flow of the boundary modes. \emph{Provenance:} \cite[Prop.~2.1]{ButlerT7} (proved self-contained there); the normalisation \cite[App.~A]{ButlerT5}; the continuation and its $2\pi$ conjugation period \cite[Prop.~7.21(1)]{ButlerT5}; the exclusion of $i\sigma$ from the real algebra \cite[Prop.~7.21(3)]{ButlerT5}.
